\documentclass[10pt]{article}

% Specify the margins and text size.
\setlength{\textwidth}{6.6in}
\setlength{\textheight}{9.5in}
\setlength{\oddsidemargin}{0pt}
\setlength{\evensidemargin}{0pt}
\setlength{\topmargin}{0pt}
\setlength{\hoffset}{.05in}
\setlength{\voffset}{-1in}

\setlength{\parskip}{5pt}
\setlength{\parindent}{0pt}

% Load some fonts and symbol packages
\usepackage{latexsym}
\usepackage{pifont}       % contains 'star' symbol for counterinsurgency handbook title
\usepackage{yfonts} 
\usepackage{amsmath}
\usepackage{amsfonts}

\usepackage{graphicx}     % actually, this is loaded by pstricks
\usepackage[T1]{fontenc}
\usepackage{ifthen}
\usepackage{pstricks,pst-grad,pst-text,pst-node,multido,pst-plot,calc,pst-3dplot}
%\usepackage[all]{xy}
%\usepackage{animate}

% The hyperref package inserts links into the pdf.
\definecolor{MyLinkColor}{rgb}{.1,.2,1}
\definecolor{MyCiteColor}{rgb}{.1,1,.2}
\definecolor{MyURLColor}{rgb}{.4,.4,.4}
\usepackage[backref=true,pagebackref=false,hyperindex,colorlinks=true,
  linkcolor=MyLinkColor,urlcolor=MyURLColor]{hyperref}


% The tweaklist package is something I found on the web.  It provides a simple interface
% for making changes to spacing used in the itemize and enumerate environments.  Comment
% this out if you don't care to use tweaklist.
\usepackage{tweaklist}
\renewcommand{\itemhook}{\setlength{\parskip}{2pt}\setlength{\parsep}%
{1pt}\setlength{\topsep}{0pt}\setlength{\itemsep}{0pt}}

\newcommand{\U}{\underline{\hspace{5pt}}}

\usepackage{listings}
\newcounter{EX}\setcounter{EX}{1}
\newcommand{\EXERCISE}{\arabic{EX}.\stepcounter{EX} }

\begin{document}
\pagestyle{empty}
\lstset{language=R, showspaces=false, showstringspaces=false}

\href{http://www.su.edu}{\includegraphics[height=1.75cm]{sulogo.eps}}
\vspace{-1.69cm}

{{\ }\hfill\small
\begin{tabular}{cl}
& Math 207\\
& Introduction to Statistics\\
\end{tabular}
}
\setlength{\baselineskip}{1.05\baselineskip}

\begin{center}
\textbf{\large  Quiz:  The Normal Curve}
\end{center}
\newcommand{\Z}{\hphantom{0}}

1. Sketch then, without using a calculator, estimate the specified area under the normal curve.

\begin{center}
\begin{tabular}{lcl}
a) $z>2$ & \hspace{3in} & b) $z>2$\\[1in]
c) $-1<z<2$ &           & d) $1<z<2$\\[1in]
e) $z<10$   &           & f) $z<-7$\hphantom{\hspace{1in}}\\[1in]
g) $-3<z<1$ &           & h) $|z|<1$\\[1in]
\end{tabular}
\end{center}

2. Use the \texttt{pnorm} function in \texttt{R} to find the specified areas under the normal curve.
Write down the commands that you use.
\begin{center}
\begin{tabular}{lcl}
a) $z<-1.5$ & \hspace{3in} & b) $z>1.8$\\[.7in]
c) $z<7.75$   &           & d) $z<-5.3$\hphantom{\hspace{1in}}\\[.7in]
e) $-1.25<z<2.25$ &           & f) $1.8<z<2.0$\\[.7in]
\end{tabular}
\end{center}

\vfill
\eject

3. Suppose that math SAT scores average out to 520 with a standard deviation of 20.
Assume that the histogram of scores follows a normal distribution.

\hspace{10pt} a) About what percent of scores were between 500 and 540?
\vspace{.8in}

\hspace{10pt} b) About what percent of scores were below 560?
\vspace{.8in}

\hspace{10pt} c) About what percent of scores were over 580?
\vspace{.8in}

\hspace{10pt} d) What score is better than 90\% of all scores? (Hint:  \texttt{qnorm}.)
\vspace{1in}

4. Errors in the measurement of a certain check weight average out to $0.0065$ kg plus or minus
about $0.0002$~kg.  Assume that the histogram of measurement errors follows a normal 
distribution.

\hspace{10pt} a) What is the chance of getting an error measurement greater than 
$0.0068$?
\vspace{.8in}

\hspace{10pt} b) What is the chance of getting an error measurement greater than 
between $0.0064$ and $0.0066$?
\vspace{.8in}

\hspace{10pt} c) About what percent of measurements  were over $0.0070$?
\vspace{.8in}

\hspace{10pt} d) What measurement is larger than 90\% of all measurements? (Hint:  \texttt{qnorm}.)
\vfill
\eject

\end{document}



